\clearpage
\appendix

\newcommand{\annexepage}[1]{%
  \clearpage
  \refstepcounter{chapter}
  \phantomsection
  \addcontentsline{toc}{chapter}{Annexe \thechapter{} : #1}
  \markboth{Annexe \thechapter{} — #1}{Annexe \thechapter{} — #1}
  \thispagestyle{empty}
  \vspace*{3.2cm}
  \begin{center}
    {\LARGE\bfseries Annexe \thechapter\par}
    \vspace{1cm}
    {\color{paulGold}\rule{\textwidth}{1pt}}\\[1.2cm]
    {\Huge\bfseries #1\par}
    \vspace{1.2cm}
    {\color{paulGold}\rule{\textwidth}{1pt}}
  \end{center}
  \vfill
  \begin{center}\color{gray}\NomProjet\end{center}
  \clearpage}

\annexepage{Guide d'installation et d'exploitation}

\section{Installation locale}

Le code et les données doivent être séparés. Le code peut être cloné depuis le
dépôt, tandis que les dossiers sensibles sont transférés directement sur la
machine de l'entreprise.

\begin{lstlisting}[style=paulpython,language=bash,caption={Commandes d'installation}]
git clone https://github.com/OthmaneW37/PrevisionPaul.git
cd PrevisionPaul
python -m pip install -r requirements.txt
python -m pytest tests -q
\end{lstlisting}

Les dossiers \texttt{data/} et \texttt{donnees\_ventes/} doivent ensuite être
placés dans la racine du projet. Ils ne sont pas publiés avec le code lorsqu'ils
contiennent des informations métier.

\section{Commandes d'exploitation}

\small
\begin{longtable}{p{6.8cm}p{7.4cm}}
\caption{Principales commandes d'exploitation}\label{tab:commandes-exploitation}\\
\toprule
\textbf{Commande} & \textbf{Effet} \\
\midrule
\endfirsthead
\toprule
\textbf{Commande} & \textbf{Effet} \\
\midrule
\endhead
\path{python dashboard.py} & Lance le tableau de bord sur \url{http://127.0.0.1:8050}. \\
\path{python main.py} & Recalcule les prévisions mensuelles, le plan sécurisé et le MRP. \\
\path{python -m paul_forecast.forecast_journalier} & Régénère l'horizon journalier de 62 jours. \\
\path{python outils/exporter_ventes_sql.py} & Lit les jours clôturés depuis SQL Server et fusionne le CSV. \\
\path{python outils/mise_a_jour_quotidienne.py} & Exécute export SQL, journalier puis mensuel avec journalisation. \\
\path{python outils/generer_tableau_recettes.py} & Crée le classeur à compléter par le chef. \\
\path{python outils/importer_recettes_chefs.py} & Importe les recettes explicitement validées. \\
\path{python -m pytest tests -q} & Lance l'ensemble des tests automatisés. \\
\bottomrule
\end{longtable}
\normalsize

\section{Procédure de mise à jour quotidienne}

\begin{enumerate}
  \item vérifier que la base Elyx est accessible ou déposer le nouvel export dans
  \path{donnees_ventes/} ;
  \item lancer la chaîne quotidienne ou attendre la tâche de 05 h 30 ;
  \item consulter le journal du jour dans \texttt{logs/} ;
  \item ouvrir le dashboard et contrôler la date « ventes réelles à jour au » ;
  \item examiner les alertes de niveau et la fiabilité avant d'exporter le plan ;
  \item vérifier le bon de commande avec le chef lorsque des recettes sont encore
  estimées.
\end{enumerate}

\annexepage{Structure du projet et fichiers produits}

\section{Arborescence fonctionnelle}

\begin{lstlisting}[basicstyle=\ttfamily\small,frame=single,backgroundcolor=\color{softGray}]
PrevisionPaul/
|-- dashboard.py              # interface Dash
|-- main.py                   # point d'entree mensuel
|-- watcher.py                # surveillance des fichiers
|-- paul_forecast/            # moteurs et regles metier
|-- data/                     # configuration JSON locale
|-- donnees_ventes/           # historique des ventes
|-- exports/                  # sorties journalieres et datees
|-- outils/                   # imports, benchmarks, recettes
|-- tests/                    # suite Pytest
|-- docs/                     # documents pour les equipes
`-- Rapport/                  # sources LaTeX du present rapport
\end{lstlisting}

\section{Exports principaux}

\small
\begin{longtable}{p{7.0cm}p{7.2cm}}
\caption{Fichiers produits par le système}\label{tab:exports}\\
\toprule
\textbf{Fichier} & \textbf{Contenu} \\
\midrule
\endfirsthead
\toprule
\textbf{Fichier} & \textbf{Contenu} \\
\midrule
\endhead
\path{exports/previsions_journalieres.csv} & Date, produit, famille, prévision, recommandation, fiabilité et commande. \\
\path{exports/suivi_prevu_reel.csv} & Comparaison hors échantillon sur les jours récents. \\
\path{plan_production_securise.csv} & Plan mensuel par produit, modèle, erreur et intervalle. \\
\path{besoins_ingredients_planifies.csv} & Quantité requise par ingrédient et par mois. \\
\path{previsions_detaillees_par_produit.csv} & Historique et scénarios de prévision détaillés. \\
\path{previsions_par_famille.csv} & Synthèse mensuelle par famille. \\
\path{validation_mensuelle_quantite.csv} & Réel et prévisions walk-forward en quantité. \\
\path{validation_mensuelle_ca.csv} & Réel et prévisions walk-forward en chiffre d'affaires. \\
\path{previsions_PAUL_2026.xlsx} & Classeur multi-feuilles à destination du responsable. \\
\path{bon_de_commande_gerant.txt} & Synthèse textuelle des matières à commander. \\
\bottomrule
\end{longtable}
\normalsize

\annexepage{Compléments de validation}

\section{Répartition des modèles sélectionnés}

Le benchmark disponible contient 848 produits. La répartition du meilleur modèle
montre l'intérêt d'une sélection hétérogène : Theta est retenu pour 293 produits,
le GBM pour 227, le naïf saisonnier pour 102, ETS pour 77, la moyenne mobile pour
44, l'ensemble médian pour 35, la moyenne pour 23, ARIMA pour 20, Croston pour 14
et l'ensemble moyen pour 13.

\begin{table}[H]
\centering
\caption{Répartition des modèles retenus dans le benchmark}
\label{tab:repartition-modeles}
\begin{tabular}{lrr@{\hspace{1.2cm}}lrr}
\toprule
\textbf{Modèle} & \textbf{Produits} & \textbf{Part} & \textbf{Modèle} & \textbf{Produits} & \textbf{Part} \\
\midrule
Theta & 293 & 34,6\,\% & Moyenne mobile & 44 & 5,2\,\% \\
GBM & 227 & 26,8\,\% & Ensemble médian & 35 & 4,1\,\% \\
Naïf saisonnier & 102 & 12,0\,\% & Moyenne & 23 & 2,7\,\% \\
ETS & 77 & 9,1\,\% & ARIMA & 20 & 2,4\,\% \\
Croston & 14 & 1,7\,\% & Ensemble moyen & 13 & 1,5\,\% \\
\bottomrule
\end{tabular}
\end{table}

\section{Graphiques statiques générés}

\begin{figure}[H]
\centering
\includegraphics[width=0.92\textwidth]{dashboard_gerant.png}
\caption{Dashboard statique destiné au responsable}
\label{fig:annexe-dashboard-gerant}
\end{figure}

\begin{figure}[H]
\centering
\includegraphics[width=0.92\textwidth]{dashboard_matieres_premieres.png}
\caption{Dashboard statique des matières premières}
\label{fig:annexe-dashboard-matieres}
\end{figure}

\section{Limites à contrôler avant usage opérationnel}

\begin{itemize}
  \item actualiser les ventes et vérifier la date de fraîcheur ;
  \item confirmer les recettes à plus fort poids matière ;
  \item remplacer les prix estimés par les factures fournisseurs ;
  \item retirer le calibrage provisoire des farines après validation complète ;
  \item analyser le biais journalier et recalibrer les seuils si nécessaire ;
  \item ne pas interpréter les marges matière comme une marge nette ;
  \item contrôler les événements et commandes futurs avant chaque recalcul.
\end{itemize}
